\chapter{TINJAUAN PUSTAKA}

\section{Landasan Teoritis}
\subsection{Pemasaran}
Pemasaran merupakan konsep dasar yang perlu dijelaskan terlebih dahulu karena penelitian ini membahas cara Rumah Makan Nasi Gerilya menyusun penawaran dan menjangkau konsumen melalui GrabFood. Pemasaran tidak hanya berarti kegiatan menjual atau mempromosikan produk, tetapi mencakup proses menciptakan, mengomunikasikan, menyampaikan, dan mempertukarkan nilai kepada pelanggan secara menguntungkan \parencite{kotler2021marketingmanagement,kotler2023principlesmarketing}. Dalam pandangan pemasaran modern, nilai pelanggan terbentuk ketika kebutuhan konsumen, kemampuan usaha, dan penawaran pasar bertemu dalam bentuk produk, harga, akses pembelian, komunikasi, serta pengalaman layanan yang dianggap relevan oleh konsumen \parencite{kotler2023principlesmarketing,wirtz2024servicesmarketing}. Karena itu, pemasaran menjadi landasan untuk memahami mengapa sebuah usaha kuliner tidak cukup hanya memiliki produk yang baik, tetapi juga perlu mengelola cara produk tersebut dipersepsikan, ditemukan, dipilih, dibeli, dan dinilai oleh konsumen.

Pada usaha kuliner, pemasaran berperan sebagai penghubung antara kemampuan internal usaha dengan kebutuhan pasar yang terus berubah. Konsumen makanan tidak hanya menilai cita rasa, tetapi juga variasi menu, keterjangkauan harga, kemudahan akses, kecepatan layanan, tampilan produk, reputasi, serta konsistensi pengalaman pembelian \parencite{regina2025kuliner,sophia2025rendanggadih,wirtz2024servicesmarketing}. Ketika pembelian dilakukan melalui platform digital, ruang pemasaran semakin meluas karena konsumen mengambil keputusan berdasarkan informasi yang tersaji di aplikasi, seperti nama menu, foto produk, deskripsi, harga, promo, rating, dan ulasan \parencite{hidayah2021digitalplatforms,putri2024cikarang,elverda2025ofdindonesia}. Dengan demikian, pemasaran dalam penelitian ini menjadi pijakan awal untuk membaca aktivitas Nasi Gerilya sebagai upaya menyusun nilai produk dan pengalaman pembelian yang sesuai dengan karakter pasar digital.

\subsection{Strategi Pemasaran}
Strategi pemasaran merupakan kelanjutan dari konsep pemasaran karena strategi menjelaskan arah dan cara usaha mengelola nilai yang ditawarkan kepada pasar sasaran. Dalam pemasaran modern, strategi tidak hanya dipahami sebagai kegiatan promosi, tetapi sebagai keputusan terpadu mengenai siapa yang dilayani, nilai apa yang ditawarkan, dan bagaimana nilai tersebut disampaikan secara lebih unggul daripada pesaing \parencite{kotler2021marketingmanagement,kotler2023principlesmarketing}. Pemahaman ini sejalan dengan pandangan manajemen strategis yang menempatkan strategi sebagai proses penyesuaian antara tujuan organisasi, kemampuan internal, dan tekanan lingkungan eksternal \parencite{david2023strategicmanagement}. Dengan demikian, strategi pemasaran perlu dibaca sebagai jembatan antara tujuan usaha dan kondisi pasar yang dihadapi, bukan sebagai kegiatan yang berdiri sendiri.

Pada usaha kuliner, strategi pemasaran semakin penting karena persaingan tidak hanya berlangsung pada kualitas rasa, tetapi juga pada variasi menu, harga, kemasan, promosi, akses pemesanan, dan pengalaman layanan yang melekat pada proses pembelian \parencite{regina2025kuliner,sophia2025rendanggadih}. Perubahan menuju pemasaran digital membuat keputusan strategis menjadi lebih kompleks karena usaha perlu mengelola tampilan produk, informasi menu, promosi, dan kanal transaksi secara bersamaan \parencite{chaffey2022digitalmarketing,hidayah2021digitalplatforms}. Dalam konteks UMKM kuliner, penelitian sebelumnya menunjukkan bahwa penguatan strategi digital dapat membantu meningkatkan visibilitas, memperluas jangkauan pasar, dan memperbaiki kinerja usaha ketika didukung oleh kesiapan organisasi serta pemilihan kanal yang tepat \parencite{putro2025digitaltransformation,regina2025kuliner}. Oleh karena itu, strategi pemasaran pada Rumah Makan Nasi Gerilya perlu dipahami sebagai rangkaian keputusan yang menghubungkan kekuatan produk, pengelolaan platform GrabFood, dan respons terhadap persaingan pasar.

Strategi tersebut harus mampu menerjemahkan pemahaman tentang konsumen menjadi keputusan yang lebih operasional, mulai dari penentuan pasar sasaran, posisi usaha, susunan menu, harga, promosi, saluran pemasaran, hingga cara usaha merespons peluang dan ancaman di dalam platform \parencite{kotler2023principlesmarketing,chaffey2022digitalmarketing,david2023strategicmanagement}. Karena fokus penelitian ini adalah perumusan strategi, konsep strategi pemasaran menjadi dasar untuk menempatkan STP, bauran pemasaran, pemasaran digital, saluran OFDA, analisis lingkungan, dan SWOT sebagai bagian yang saling melengkapi.

\subsection{STP dan Bauran Pemasaran}
Segmentasi, \textit{targeting}, dan \textit{positioning} (STP) melengkapi strategi pemasaran karena membantu usaha menentukan kelompok pasar yang ingin dilayani dan citra yang ingin dibangun di benak konsumen \parencite{kotler2021marketingmanagement,kotler2023principlesmarketing}. Segmentasi diperlukan karena pasar kuliner tidak homogen; konsumen dapat berbeda dalam daya beli, waktu pembelian, preferensi menu, sensitivitas harga, dan respons terhadap promosi \parencite{suharjo2020bogorsnacks,ruhiyat2023minatbeli}. Setelah segmen dipahami, usaha perlu memilih segmen yang paling sesuai dengan kapasitas dan tujuan usahanya, lalu membangun \textit{positioning} yang membedakan penawarannya dari pesaing \parencite{kotler2023principlesmarketing,sophia2025rendanggadih}. Pada konteks Nasi Gerilya, STP menjadi penting karena strategi pemasaran di GrabFood perlu menjelaskan pasar sasaran, posisi harga, identitas rasa, dan alasan konsumen memilih usaha tersebut dibanding \textit{merchant} sejenis.

Bauran pemasaran menerjemahkan STP menjadi tindakan pemasaran yang lebih operasional. Konsep 4P menempatkan \textit{product}, \textit{price}, \textit{place}, dan \textit{promotion} sebagai perangkat utama untuk mengelola penawaran pasar \parencite{kotler2021marketingmanagement,kotler2023principlesmarketing}. Dalam usaha kuliner digital, produk tidak hanya berupa makanan, tetapi juga variasi menu, kemasan, foto, dan kejelasan deskripsi; harga tidak hanya berupa nominal menu, tetapi juga potongan harga, ongkos kirim, dan paket promosi; tempat tidak hanya berarti lokasi fisik, tetapi juga kanal distribusi digital; sedangkan promosi mencakup voucher, kampanye platform, dan komunikasi visual \parencite{chaffey2022digitalmarketing,hidayah2021digitalplatforms,putri2024cikarang}. Penelitian tentang strategi pemasaran kuliner berbasis platform juga menunjukkan bahwa penerapan STP dan bauran pemasaran dapat membantu usaha mengarahkan penjualan ketika fitur platform dimanfaatkan secara konsisten \parencite{sophia2025rendanggadih,firdaus2022bangkupawon}.

Pada bisnis jasa dan kuliner, bauran pemasaran sering diperluas menjadi 7P dengan menambahkan \textit{people}, \textit{process}, dan \textit{physical evidence}, karena keberhasilan pemasaran tidak hanya ditentukan oleh penawaran inti, tetapi juga oleh pelaksana layanan, proses pemenuhan pesanan, dan bukti yang memperkuat persepsi kualitas \parencite{wirtz2024servicesmarketing}. Dalam konteks OFDA, unsur tersebut terlihat dari kesiapan admin dan karyawan, kecepatan konfirmasi, akurasi pesanan, pengelolaan menu habis, tampilan etalase digital, foto produk, kemasan, rating, dan ulasan \parencite{hidayah2021digitalplatforms,putri2024cikarang,elverda2025ofdindonesia}. Dengan demikian, STP dan bauran pemasaran saling melengkapi: STP membantu menentukan arah pasar dan posisi usaha, sedangkan bauran pemasaran membantu menilai apakah keputusan produk, harga, distribusi, promosi, proses, pelaksana layanan, dan bukti digital sudah mendukung posisi tersebut.

\subsection{Pemasaran Digital}
Pemasaran digital memperluas pembahasan strategi pemasaran karena teknologi digital mengubah cara usaha menciptakan, mengomunikasikan, dan menyampaikan nilai kepada pasar \parencite{chaffey2022digitalmarketing}. Berbeda dari pemasaran konvensional yang sering bergantung pada komunikasi satu arah, pemasaran digital memungkinkan interaksi, pengukuran respons, dan penyesuaian pesan secara lebih cepat melalui data kunjungan, klik, ulasan, transaksi, dan keterlibatan konsumen \parencite{chaffey2022digitalmarketing,purnomo2023conversion}. Dalam konteks usaha kuliner, pemasaran digital tidak hanya berarti mengunggah konten promosi, tetapi juga mengelola katalog menu, foto produk, deskripsi, promosi, kanal pemesanan, dan reputasi digital \parencite{regina2025kuliner,hidayah2021digitalplatforms}. Dengan demikian, pemasaran digital menjadi bagian dari strategi bisnis yang menghubungkan aktivitas promosi, distribusi, transaksi, dan evaluasi pasar.

Pemasaran digital relevan bagi UMKM kuliner karena usaha kecil sering menghadapi keterbatasan lokasi, modal promosi, dan jangkauan pelanggan, sedangkan platform digital membuka peluang untuk memperluas pasar dengan biaya yang relatif lebih fleksibel \parencite{fatmawatie2022turnover,regina2025kuliner}. Penelitian \textcite{putro2025digitaltransformation} menunjukkan bahwa transformasi digital berpengaruh terhadap kinerja UMKM kuliner melalui dukungan teknologi, organisasi, dan lingkungan. Kajian lain menunjukkan bahwa penggunaan media sosial, katalog daring, \textit{e-commerce}, dan aplikasi pesan antar dapat meningkatkan visibilitas, penjualan, dan keterhubungan usaha dengan konsumen ketika dikelola secara konsisten \parencite{hidayah2021digitalplatforms,indriyani2025gudangbaju,pratama2024sambelmoji}. Oleh karena itu, pemasaran digital dalam penelitian ini diposisikan sebagai konteks strategis yang membantu membaca bagaimana Nasi Gerilya memanfaatkan GrabFood sebagai kanal pemasaran, bukan sekadar sebagai aplikasi pemesanan.

\subsection{Saluran Pemasaran (OFDA)}
\textit{Online food delivery application} (OFDA) merupakan saluran digital yang mempertemukan usaha makanan dengan konsumen melalui sistem yang mengintegrasikan informasi produk, pemesanan, pembayaran, promosi, dan pengantaran \parencite{zaheer2024ofda,elverda2025ofdindonesia}. Dalam kajian strategi pemasaran, OFDA tidak hanya dipahami sebagai alat distribusi, tetapi juga sebagai etalase digital, ruang promosi, sumber reputasi, dan arena persaingan antar \textit{merchant} \parencite{chaffey2022digitalmarketing,putri2024cikarang}. Karakteristik ini membuat keputusan strategi pada platform seperti GrabFood perlu memperhatikan keterpaduan antara menu, harga, foto, deskripsi, promo, ketersediaan toko, rating, dan posisi terhadap pesaing \parencite{hidayah2021digitalplatforms,firdaus2022bangkupawon}.

Peran OFDA dalam usaha kuliner berkaitan langsung dengan bauran pemasaran karena produk ditampilkan melalui nama menu dan visual, harga ditampilkan bersama potongan dan ongkos kirim, distribusi dilakukan melalui jaringan pengantaran, sedangkan promosi dijalankan melalui voucher, kampanye platform, dan paket bundling \parencite{putri2024cikarang,sophia2025rendanggadih}. Penelitian tentang pemanfaatan GrabFood dan GoFood menunjukkan bahwa platform pesan antar dapat membantu UMKM memperluas jangkauan pasar dan meningkatkan penjualan, tetapi manfaat tersebut tetap bergantung pada kemampuan usaha mengelola fitur platform secara tepat \parencite{firdaus2022bangkupawon,putri2024cikarang}. Dengan demikian, OFDA dalam penelitian ini menjadi konteks strategis untuk membaca bagaimana Nasi Gerilya mengelola saluran GrabFood dan bagaimana dinamika platform tersebut memunculkan peluang maupun ancaman bagi strategi pemasaran.

\subsection{Analisis Lingkungan}
Analisis lingkungan diperlukan karena strategi pemasaran harus disusun berdasarkan kondisi nyata usaha dan dinamika pasar yang dihadapi. Dalam manajemen strategis, lingkungan internal menggambarkan sumber daya dan kemampuan yang berada dalam kendali usaha, sedangkan lingkungan eksternal menggambarkan kondisi pasar, pesaing, teknologi, kebijakan, dan perubahan perilaku konsumen yang memengaruhi ruang gerak usaha \parencite{david2023strategicmanagement,kotler2023principlesmarketing}. Dalam kerangka SWOT, faktor internal menjadi dasar untuk menentukan kekuatan dan kelemahan, sedangkan faktor eksternal menjadi dasar untuk menentukan peluang dan ancaman \parencite{gurel2017swot,helms2010swotreview}. Dengan demikian, analisis lingkungan membantu penelitian ini menghindari perumusan strategi yang hanya bertumpu pada keinginan usaha, karena strategi perlu disesuaikan dengan kapasitas internal dan tekanan eksternal.

\subsubsection{Lingkungan Internal}
Lingkungan internal pada usaha kuliner dapat dilihat melalui beberapa bidang, seperti sumber daya manusia, keuangan, produksi, operasional, pemasaran, dan teknologi. Aspek sumber daya manusia berkaitan dengan jumlah karyawan, pembagian tugas, kemampuan mengelola pesanan, kecepatan pelayanan, dan kedisiplinan dalam menjaga standar kerja \parencite{wirtz2024servicesmarketing,david2023strategicmanagement}. Aspek keuangan berkaitan dengan kemampuan usaha mengatur harga, biaya bahan baku, biaya promosi, potongan platform, dan margin penjualan agar strategi promosi tidak menurunkan keberlanjutan usaha \parencite{kotler2023principlesmarketing,david2023strategicmanagement}. Aspek produksi dan operasional mencakup konsistensi rasa, kualitas bahan, variasi menu, ketersediaan stok, pengemasan, akurasi pesanan, dan ketepatan waktu pemrosesan pesanan \parencite{regina2025kuliner,sophia2025rendanggadih,putri2024cikarang}. Aspek pemasaran dan teknologi mencakup kemampuan usaha memperbarui menu, mengelola foto dan deskripsi produk, menjalankan promo, membaca ulasan, serta memanfaatkan fitur GrabFood sebagai kanal pemasaran \parencite{chaffey2022digitalmarketing,hidayah2021digitalplatforms,firdaus2022bangkupawon}. Keseluruhan aspek tersebut menentukan apakah kekuatan Nasi Gerilya dapat dimaksimalkan dan apakah kelemahannya perlu diperbaiki sebelum strategi baru dijalankan.

\subsubsection{Lingkungan Eksternal}
Lingkungan eksternal melengkapi analisis internal karena usaha kuliner berbasis platform tidak beroperasi dalam ruang yang sepenuhnya dapat dikendalikan. Faktor eksternal dapat mencakup pertumbuhan penggunaan platform pesan antar, kebiasaan konsumen mencari makanan secara daring, kebijakan biaya layanan, program promo platform, ongkos kirim, algoritma visibilitas toko, rating, ulasan, serta kehadiran \textit{merchant} pesaing dengan produk dan harga yang sebanding \parencite{hidayah2021digitalplatforms,putri2024cikarang,elverda2025ofdindonesia}. Penelitian tentang strategi digital UMKM menunjukkan bahwa media digital dan platform dapat membuka peluang peningkatan penjualan, tetapi peluang tersebut dapat melemah ketika usaha menghadapi tekanan pesaing, perang promo, keterbatasan modal promosi, atau ketidakmampuan mengelola kanal digital secara konsisten \parencite{saptaji2025karawang,azzahra2024guzzini,indriyani2025gudangbaju}. Oleh karena itu, analisis lingkungan pada Nasi Gerilya digunakan untuk membaca hubungan antara kapasitas internal seperti SDM, keuangan, operasional, dan pemasaran dengan kondisi eksternal seperti kebijakan GrabFood, perilaku konsumen, persaingan kategori, serta perubahan saluran pemasaran digital.

\subsection{SWOT}
Analisis SWOT menjadi alat utama dalam penelitian ini karena SWOT menghubungkan faktor internal dan faktor eksternal ke dalam empat kategori strategis, yaitu \textit{strengths}, \textit{weaknesses}, \textit{opportunities}, dan \textit{threats} \parencite{helms2010swotreview,gurel2017swot}. Kekuatan dan kelemahan berasal dari kondisi internal usaha, sedangkan peluang dan ancaman berasal dari lingkungan eksternal yang memengaruhi arah keputusan strategis \parencite{david2023strategicmanagement,gurel2017swot}. Keunggulan SWOT terletak pada kemampuannya menyusun informasi yang tersebar menjadi peta strategis yang mudah dibaca, tetapi kajian terdahulu juga menekankan bahwa SWOT perlu digunakan secara sistematis agar tidak berhenti sebagai daftar faktor semata \parencite{helms2010swotreview,david2023strategicmanagement}. Oleh karena itu, SWOT dalam penelitian ini digunakan bukan hanya untuk mengidentifikasi kondisi usaha, tetapi juga untuk menghubungkan temuan lapangan dengan alternatif strategi pemasaran.

Matriks SWOT memperdalam analisis karena faktor internal dan eksternal tidak hanya diklasifikasikan, tetapi dikombinasikan menjadi strategi SO, WO, ST, dan WT. Strategi SO memanfaatkan kekuatan untuk mengambil peluang, strategi WO menggunakan peluang untuk mengatasi kelemahan, strategi ST memakai kekuatan untuk menghadapi ancaman, sedangkan strategi WT diarahkan untuk meminimalkan kelemahan dan menghindari ancaman \parencite{david2023strategicmanagement,gurel2017swot}. Dalam penelitian UMKM dan usaha kuliner, penggunaan SWOT telah banyak dipakai untuk merumuskan strategi pemasaran karena mampu mempertemukan persoalan produk, promosi, kapasitas internal, dan dinamika persaingan pasar \parencite{azzahra2024guzzini,saptaji2025karawang,haliza2023seqapla}. Dengan demikian, matriks SWOT relevan untuk membantu merumuskan strategi pemasaran Nasi Gerilya secara lebih operasional berdasarkan data observasi, wawancara, ulasan, dan dokumentasi performa.

Alur teori dalam penelitian ini berakhir pada SWOT karena seluruh konsep sebelumnya saling memberi dasar untuk proses perumusan strategi. Pemasaran menjelaskan penciptaan nilai bagi pelanggan, strategi pemasaran menjelaskan arah persaingan usaha, STP dan bauran pemasaran menjelaskan penerjemahan strategi ke pasar sasaran dan unsur operasional, pemasaran digital serta OFDA menjelaskan konteks GrabFood sebagai saluran pemasaran, sedangkan analisis lingkungan membantu memilah faktor SDM, keuangan, operasional, pemasaran, teknologi, konsumen, pesaing, dan kebijakan platform \parencite{kotler2023principlesmarketing,chaffey2022digitalmarketing,david2023strategicmanagement}. Seluruh rangkaian tersebut kemudian digunakan untuk mengidentifikasi kekuatan, kelemahan, peluang, dan ancaman Rumah Makan Nasi Gerilya sebelum dirumuskan menjadi alternatif strategi pemasaran melalui matriks SWOT \parencite{helms2010swotreview,gurel2017swot,saptaji2025karawang}.

\clearpage
\begin{landscape}
    \section{Penelitian Terdahulu}

    Penelitian terdahulu perlu ditelaah untuk melihat posisi penelitian ini di antara kajian yang membahas strategi pemasaran digital, pemanfaatan platform \textit{online food delivery}, dan perumusan strategi berbasis analisis SWOT \parencite{regina2025kuliner,putri2024cikarang,saptaji2025karawang,azzahra2024guzzini}. Ringkasan penelitian terdahulu yang relevan disajikan pada Tabel~\ref{tab:penelitian-terdahulu}.

    {\footnotesize
    \setlength{\tabcolsep}{4pt}
    \renewcommand{\arraystretch}{1.10}
    \setlength{\LTleft}{0pt}
    \setlength{\LTright}{0pt}
    \setlength{\LTcapwidth}{\linewidth}
    \newlength{\PenelitianTerdahuluWidth}
    \setlength{\PenelitianTerdahuluWidth}{\dimexpr\linewidth-6\tabcolsep-3pt\relax}
    \refstepcounter{table}\label{tab:penelitian-terdahulu}
    \edef\PenelitianTerdahuluNomor{\thetable}
    \addcontentsline{lot}{table}{\protect\numberline{\thetable}Penelitian Terdahulu yang Relevan}
    \begin{longtable}{@{}>{\raggedright\arraybackslash}p{0.14\PenelitianTerdahuluWidth}>{\raggedright\arraybackslash}p{0.30\PenelitianTerdahuluWidth}>{\raggedright\arraybackslash}p{0.18\PenelitianTerdahuluWidth}>{\raggedright\arraybackslash}p{0.38\PenelitianTerdahuluWidth}@{}}
        \multicolumn{4}{@{}c@{}}{\textbf{Tabel \PenelitianTerdahuluNomor}}\\
        \multicolumn{4}{@{}c@{}}{\textbf{Penelitian Terdahulu yang Relevan}}\\[6pt]
        \toprule
        \textbf{Nama (Tahun)}                     & \textbf{Judul Penelitian}                                                                                                                                                            & \textbf{Metode Analisis}                                                                                                                       & \textbf{Hasil Penelitian}                                                                                                                                                                                                                                           \\
        \midrule
        \endfirsthead
        \multicolumn{4}{@{}l@{}}{\textbf{Lanjutan Tabel \PenelitianTerdahuluNomor}}\\
        \toprule
        \textbf{Nama (Tahun)}                     & \textbf{Judul Penelitian}                                                                                                                                                            & \textbf{Metode Analisis}                                                                                                                       & \textbf{Hasil Penelitian}                                                                                                                                                                                                                                           \\
        \midrule
        \endhead
        \midrule
        \multicolumn{4}{r}{\footnotesize Dilanjutkan pada halaman berikutnya}\\
        \endfoot
        \bottomrule
        \endlastfoot
        \textcite{regina2025kuliner}              & Strategi Pemasaran Digital untuk Meningkatkan Penjualan UMKM Sektor Kuliner                                                                                                          & Deskriptif kualitatif                                                                                                               & Strategi digital melalui media sosial, katalog daring, dan pemesanan online meningkatkan visibilitas usaha, interaksi dengan konsumen, dan peluang kenaikan penjualan UMKM kuliner.                                                      \\
        \textcite{sophia2025rendanggadih}         & Analisis Strategi Pemasaran Rendang Gadih Melalui Platform Shopee untuk Meningkatkan Penjualan                                                                                       & Deskriptif kualitatif                                                                                                               & Peningkatan penjualan dicapai melalui penerapan STP, bauran pemasaran 4P, serta pemanfaatan fitur Shopee seperti CRM, Live, Ads, dan Campaign.                                                   \\
        \textcite{hendrizon2024madayacoffee}      & Analisis Strategi Pemasaran pada Madaya Coffee Parung Bogor                                                                                                                          & Matriks IFE, EFE, SWOT, dan QSPM                                                                                                    & Madaya Coffee perlu meningkatkan aktivitas mengikuti pameran untuk memperkenalkan produk kepada khalayak lebih luas sebagai bagian dari upaya peningkatan penjualan.                              \\
        \textcite{ramadhanti2023rucika}           & Analisis SWOT dalam Menentukan Strategi Pemasaran Rucika pada Distributor CV Sinar Mas Agung Semarang                                                                                & Kualitatif dengan studi kasus; wawancara, observasi, dokumentasi, SWOT, dan bauran pemasaran 7P                         & Strategi pemasaran CV Sinar Mas Agung mengarah pada \textit{growth oriented strategy}; aspek promosi dan bukti kepuasan konsumen masih perlu dimaksimalkan.                                      \\
        \textcite{haliza2023seqapla}              & Analisis Strategi Pemasaran dalam Upaya Peningkatan Penjualan pada Seqapla Coffee                                                                                                    & Analisis deskriptif kondisi internal-eksternal dan perumusan strategi pemasaran                                                     & Seqapla Coffee membutuhkan penguatan strategi pemasaran yang lebih terarah untuk menghadapi persaingan \textit{coffee shop} dan mendorong pemulihan penjualan.                                  \\
        \textcite{ruhiyat2023minatbeli}           & Strategi Pemasaran Online untuk Meningkatkan Minat Beli Konsumen Produk Makanan dan Minuman UMKM di Kota Bogor                                                                       & Regresi linear berganda, analisis deskriptif, SWOT, IFE-EFE, IE, dan QSPM                                                           & Unsur 4P berpengaruh positif terhadap minat beli; posisi UMKM berada pada strategi \textit{grow and build} dengan rekomendasi inovasi produk sesuai kebutuhan dan daya beli konsumen.             \\
        \textcite{khairani2022buahsayur}          & Strategi Meningkatkan Penjualan Buah dan Sayur pada Platform Digital di Kota Bogor                                                                                                   & Analisis deskriptif, regresi data panel, dan SWOT                                                                                   & Harga, ulasan, peringkat, jenis toko, dan promosi memengaruhi penjualan; strategi yang direkomendasikan adalah kerja sama langsung dengan produsen sebagai pemasok utama.                        \\
        \textcite{mamuriyah2021gadihminang}       & Strategi Pemasaran dalam Meningkatkan Hasil Penjualan Rumah Makan Padang Gadih Minang                                                                                                & Kegiatan PKM dengan promosi digital melalui brosur, Instagram, dan YouTube                                                          & Promosi melalui brosur, Instagram, dan YouTube diarahkan untuk memperkenalkan Rumah Makan Padang Gadih Minang ke khalayak yang lebih luas dan diharapkan mendorong kenaikan omzet penjualan.    \\
        \textcite{roslam2021markisaaurora}        & Strategi Pemasaran UMKM Markisa Aurora di Kota Makassar pada Masa Pandemi COVID-19                                                                                                   & Analisis deskriptif strategi pemasaran UMKM pada masa pandemi                                                                       & UMKM Markisa Aurora dapat mempertahankan dan memperluas pasar melalui penguatan promosi digital, penyesuaian strategi pemasaran, dan pemanfaatan kanal penjualan daring selama pandemi.          \\
        \textcite{suharjo2020bogorsnacks}         & Strategi Pemasaran Digital pada UMKM Makanan Ringan di Kota Bogor                                                                                                                    & Analisis deskriptif, perilaku pembelian daring konsumen, \textit{Importance Performance Analysis}, bauran pemasaran digital, dan 4C & Strategi prioritas menekankan penguatan konten digital, pemanfaatan media sosial, dan penyesuaian elemen bauran pemasaran digital untuk memperluas jangkauan pasar UMKM makanan ringan.        \\
        \textcite{gurel2017swot}                  & SWOT Analysis: A Theoretical Review                                                                                                                                                  & Tinjauan teoretis terhadap konsep dan penggunaan SWOT                                                                               & Kajian ini menegaskan bahwa SWOT merupakan alat analisis strategis yang membantu organisasi mengidentifikasi kekuatan, kelemahan, peluang, dan ancaman sebagai dasar perumusan strategi.         \\
        \textcite{helms2010swotreview}            & Exploring SWOT Analysis -- Where Are We Now?: A Review of Academic Research from the Last Decade                                                                                     & Kajian pustaka terhadap penelitian SWOT dalam satu dekade                                                                           & Hasil telaah menunjukkan bahwa SWOT tetap luas digunakan dalam penelitian strategi, tetapi memerlukan penerapan yang lebih sistematis dan dukungan kerangka analitis yang lebih kuat.             \\
    \end{longtable}
    \tablesource[\linewidth]{Diolah peneliti dari berbagai penelitian terdahulu (2026).}
    \addtocounter{table}{-1}
}

\end{landscape}
\clearpage

\section{Kerangka Konseptual}
Kerangka konseptual dalam penelitian kualitatif deskriptif berfungsi sebagai alur berpikir yang membantu peneliti menjaga hubungan antara masalah penelitian, landasan teori, data lapangan, dan proses analisis \parencite{creswell2018qualitativeinquiry,miles2020qualitativedata}. Dalam penelitian ini, kerangka konseptual tidak digunakan untuk menguji hubungan sebab akibat antarvariabel, tetapi untuk mengarahkan proses identifikasi faktor internal, faktor eksternal, dan perumusan strategi pemasaran Rumah Makan Nasi Gerilya pada platform GrabFood. Alur tersebut dibangun dari konsep pemasaran, strategi pemasaran, STP, bauran pemasaran, pemasaran digital, saluran pemasaran OFDA, analisis lingkungan, dan SWOT \parencite{kotler2023principlesmarketing,chaffey2022digitalmarketing,david2023strategicmanagement,gurel2017swot}.

Berdasarkan Gambar~\ref{fig:kerangka-konseptual}, penelitian dimulai dari objek penelitian, yaitu Rumah Makan Nasi Gerilya pada platform GrabFood, kemudian diarahkan pada kondisi persaingan \textit{merchant} kuliner yang semakin ketat dan masalah performa usaha yang belum stabil. Masalah tersebut dibaca melalui landasan teori pemasaran, strategi pemasaran, STP, bauran pemasaran, pemasaran digital, dan saluran pemasaran OFDA. Data penelitian kemudian dipilah menjadi faktor internal dan faktor eksternal, karena kedua kelompok faktor tersebut menjadi dasar penentuan kekuatan, kelemahan, peluang, dan ancaman \parencite{david2023strategicmanagement,gurel2017swot}. Hasil pemetaan faktor tersebut dimasukkan ke dalam matriks SWOT untuk menyusun alternatif strategi SO, WO, ST, dan WT, lalu diarahkan menjadi rumusan strategi pemasaran Rumah Makan Nasi Gerilya pada GrabFood. Karena penelitian ini bersifat kualitatif deskriptif, penelitian tidak merumuskan hipotesis statistik; kerangka konseptual digunakan sebagai pedoman analisis untuk menjawab pertanyaan penelitian secara sistematis \parencite{creswell2018qualitativeinquiry,miles2020qualitativedata}.

\begin{center}
    \centering
    \begingroup
    \setlength{\fboxsep}{2pt}
    \setlength{\fboxrule}{0.35pt}
    \newcommand{\komponen}[2][2.35cm]{%
        \fcolorbox{black!55}{white}{%
            \parbox[c][0.74cm][c]{#1}{\centering\scriptsize\linespread{0.86}\selectfont #2}%
        }%
    }
    \newcommand{\komponenbesar}[2][5.05cm]{%
        \fcolorbox{black!55}{white}{%
            \parbox[c][0.78cm][c]{#1}{\centering\scriptsize\linespread{0.9}\selectfont #2}%
        }%
    }
    \newcommand{\grup}[2]{%
        \begin{tabular}{@{}c@{}}
            \textbf{#1} \\[2pt]
            #2
        \end{tabular}%
    }
    \newcommand{\empatkomponen}[4]{%
        \begin{tabular}{@{}c@{\hspace{2pt}}c@{\hspace{2pt}}c@{\hspace{2pt}}c@{}}
            \komponen{#1} & \komponen{#2} & \komponen{#3} & \komponen{#4}
        \end{tabular}%
    }
    \newcommand{\duakomponen}[2]{%
        \begin{tabular}{@{}c@{\hspace{4pt}}c@{}}
            \komponenbesar{#1} & \komponenbesar{#2}
        \end{tabular}%
    }
    \tikzset{
        drop shadow/.style={},
        diagram arrow type/.append style={draw=black!70}
    }
    \smartdiagramset{
        module shape=rectangle,
        module minimum width=11.7cm,
        module minimum height=0.74cm,
        module y sep=1.83,
        text width=11.05cm,
        font=\scriptsize\linespread{0.94}\selectfont,
        border color=black!65,
        arrow color=black!70,
        arrow line width=0.8pt,
        arrow tip=stealth,
        back arrow disabled=true,
        set color list={black!10,white,white,black!4,black!4,black!5,black!10,black!5,black!10}
    }
    \smartdiagram[flow diagram]{
        {\textbf{Rumah Makan Nasi Gerilya}},
        {Persaingan \textit{merchant} kuliner pada platform GrabFood semakin ketat},
        {\textbf{Masalah Penelitian}\\Performa penjualan; retensi pelanggan; daya saing kategori belum stabil},
        {\grup{Landasan Teoritis}{\empatkomponen{Pemasaran\\Strategi Pemasaran}{STP\\Bauran Pemasaran}{Pemasaran\\Digital}{Saluran\\OFDA}}},
        {\textbf{Data Lapangan}\\Observasi GrabFood; wawancara; ulasan pelanggan; pra-survei; dokumentasi performa usaha},
        {\grup{Analisis Lingkungan}{\duakomponen{\textbf{Internal}\\SDM; keuangan; operasional; produk; pemasaran; teknologi}{\textbf{Eksternal}\\Konsumen; pesaing; kebijakan platform; promo; rating; ulasan}}},
        {\textbf{Matriks SWOT}\\Pemetaan kekuatan; kelemahan; peluang; ancaman},
        {\grup{Alternatif Strategi}{\empatkomponen{\textbf{SO}\\Strength--Opportunity}{\textbf{WO}\\Weakness--Opportunity}{\textbf{ST}\\Strength--Threat}{\textbf{WT}\\Weakness--Threat}}},
        {\textbf{Rumusan Strategi Pemasaran}\\Rumah Makan Nasi Gerilya pada GrabFood}
    }
    \endgroup
    \captionof{figure}{Kerangka Konseptual Penelitian}
    \label{fig:kerangka-konseptual}
\end{center}